\documentclass{article}
\usepackage{fontspec} % XeLaTeX hanterar UTF-8 automatiskt
\usepackage{hyperref} % klickbara länkar (valfritt)
\usepackage{enumitem}

\title{Installation och Körning av Team04 Applikation}
\author{}
\date{}

\begin{document}

\maketitle

\section*{Installation och Uppstart}

\subsection*{1. Förberedelser}
\begin{itemize}[leftmargin=*]
    \item Kontrollera att servern har \textbf{Python 3.14} och \textbf{Node.js 24} installerade.
\end{itemize}

\subsection*{2. Ladda ner och packa upp}
\begin{itemize}[leftmargin=*]
    \item Hämta zip-filen som innehåller \texttt{app} mappen och dokumentation.
    \item Packa upp zip-filen på servern i önskad katalog.
\end{itemize}

\subsection*{3. Backend}
\begin{itemize}[leftmargin=*]
    \item Backend är redan byggd som en fristående binär med \texttt{pyinstaller}, så inga extra Python-bibliotek behövs.
\end{itemize}

\subsection*{4. Frontend}
\begin{itemize}[leftmargin=*]
    \item Frontend är redan byggd och ligger i \texttt{app/frontend/dist}.
    \item Inga extra byggkommandon behövs på servern.
\end{itemize}

\subsection*{5. Starta applikationen}
\begin{itemize}[leftmargin=*]
    \item Gå till \texttt{app/backend}.
    \item Kör backend-binären:
    \begin{verbatim}
    ./run
    \end{verbatim}
    \item Applikationen startar på den host och port som är satt innan programet kompilerades. (exempel: \texttt{http://100.84.14.107:8000}).
\end{itemize}

\subsection*{6. Verifiering}
\begin{itemize}[leftmargin=*]
    \item Öppna webbläsaren och navigera till serverns IP och port, eller eventuell URL.
\end{itemize}

\section*{Noteringar}
\begin{itemize}[leftmargin=*]
    \item Inga utvecklingsberoenden behövs på servern.
    \item Endast \texttt{app}-mappen behöver laddas upp.
\end{itemize}

\end{document}